\section{The implementation}
\label{sec:impl}

\noindent\emph{Status: draft -- the module map, the data layout and the API example describe the code as it
stands and have been re-checked against it, including a run of the example reproduced below; the test inventory
is quoted from a full run of the suite.}

The theory of Sections~\ref{sec:theory} and~\ref{sec:energy} occupies far fewer lines than the code that carries
it. That is expected, and it is worth saying why: the mathematics of the finite integration technique is a
handful of incidence matrices and diagonal star matrices, but an implementation has to decide where each of them
lives, how the unknowns are numbered, which quantity is a vector and which a sparse matrix, which arrays are
cached, and how the conductor constraint is expressed so that it costs nothing at run time. Those decisions are
what this section records, because they are the part of the package that a reader wanting to change it needs
first.

The package is pure Python on top of \texttt{numpy} and \texttt{scipy}. There is no compiled extension and no
generated code, and the only dependency beyond those two is \texttt{matplotlib}, which is imported by the
plotting module alone and never by the solver. The whole of \texttt{emfit} is about $6.4$ thousand lines,
distributed as in Table~\ref{tab:modules}; the ten test modules add a further $2.4$ thousand lines and the two
example scripts another four hundred, which is a fair reflection of where the effort went.

\subsection{The module map}
\label{sec:modules}

Each module owns exactly one of the objects of Section~\ref{sec:theory}, and the dependency arrows point one way.

\begin{table}[H]
  \centering
  \caption{The modules of \texttt{emfit}, their size, and the object each one owns. The last column points at
  the part of the theory that fixes the module's behaviour.}
  \label{tab:modules}
  \small
  \begin{tabular}{@{}l r p{7.0cm} p{3.6cm}@{}}
    \toprule
    Module & Lines & Owns & Governing equations \\
    \midrule
    \texttt{grid.py}       & 254  & node coordinates, widths, dual widths, areas, volumes & \eqref{eq:stars} \\
    \texttt{operators.py}  & 326  & index maps, $\Gmat$, $\Cmat$, $\Smat$, duals, PEC mask & \eqref{eq:identities}, \eqref{eq:pec} \\
    \texttt{materials.py}  & 266  & the diagonal stars $\Em$, $\Mm$, $\Ms$, constants & \eqref{eq:stars} \\
    \texttt{evolve.py}     & 276  & leapfrog step, sources, CFL bound, energy & \eqref{eq:leapfrog-h}--\eqref{eq:invariant} \\
    \texttt{maxwell.py}    & 103  & curl--curl assembly, PEC constraints & \eqref{eq:rcurl} \\
    \texttt{eigen.py}      & 261  & resonances, static cut, residual screen, analytic box & \eqref{eq:geneig-f}, \eqref{eq:cut} \\
    \texttt{abc.py}        & 348  & Mur planes, neighbour map, coefficient $K$ & \eqref{eq:mur} \\
    \texttt{pml.py}        & 461  & conductivity profile, damping multipliers, layer fit & \eqref{eq:pml-match}, \eqref{eq:pml-profile} \\
    \texttt{scatter.py}    & 565  & staircase PEC cylinder, background subtraction & --- \\
    \texttt{nearfield.py}  & 922  & DFT accumulator, contour integral, far field, RCS & \eqref{eq:sigma2d} \\
    \texttt{mie.py}        & 295  & analytic two-dimensional PEC cylinder series & --- \\
    \texttt{studies.py}    & 1182 & the named runs behind every front end & --- \\
    \texttt{viz.py}        & 352  & spectrum, field slice, energy panels & --- \\
    \texttt{webapp.py}     & 431  & routes, forms, rendered figures & --- \\
    \texttt{cli.py}        & 191  & argument parsing, table output & --- \\
    \texttt{\_env.py}      & 70   & thread pinning before \texttt{numpy} & --- \\
    \texttt{\_\_init\_\_.py} & 41 & version and submodule imports & --- \\
    \texttt{\_\_main\_\_.py} & 13 & \texttt{python -m} entry point & --- \\
    \midrule
    total                  & 6357 &                                                      & \\
    \bottomrule
  \end{tabular}
\end{table}

The split that matters is \texttt{grid} against \texttt{operators}. \texttt{grid.py} holds a frozen dataclass
whose entire content is the three coordinate arrays; it answers metric questions -- the width of a cell, the
length of an edge, the area of a face, the volume of a cell -- and it knows nothing about topology, numbering or
materials. \texttt{operators.py} holds the \texttt{CellComplex}, which answers topological questions -- how many
edges, which edge is which, what $\Cmat$ and $\Smat$ are -- and it knows nothing about lengths. The claim tested
in \texttt{test\_operators} is precisely that: the incidence matrices of a uniform grid and of a deliberately
jittered non-uniform grid of the same cell counts are the same integer arrays, because refinement can only change
the stars and never the connectivity. A package that mixed the two would still produce plausible numbers on a
uniform grid and would fail silently on a graded one.

The second split is time domain against frequency domain. \texttt{evolve.py} and \texttt{eigen.py} share the stars
and the operators and nothing else; the stepper never assembles a matrix, and the eigensolver never takes a step.
\texttt{maxwell.py} is the bridge: a thin class that assembles the curl--curl system of Section~\ref{sec:eigen}
for a caller who wants a system object rather than a function call, and applies the conductor constraint as a
restriction of the unknown space rather than as a penalty. It is the smallest module that carries physics, which is the
intended outcome of having written the operators first.

\subsection{Data layout and sparse operators}
\label{sec:layout}

Every field quantity is a flat one-dimensional \texttt{float64} array, and every operator is a
\texttt{scipy.sparse} matrix in CSR form. The two statements are related: the operators have at most six
non-zeros per row, so a matrix--vector product is a short gather, and that gather is the whole of the spatial
discretisation at each step. The number of unknowns in each family follows from the cell counts through the
relations tabulated in Table~\ref{tab:counts}; the $10\times7\times5$ cavity of the demonstration has $1382$
edges, of which $762$ remain free once the six walls are constrained, and the $40\times6\times6$ validation grid
has $5404$ edges and $4836$ faces.

Numbering within each family is a stride formula, and it is not something a caller should ever write out: the
complex exposes \texttt{idx\_ex}, \texttt{idx\_ey}, \texttt{idx\_ez}, \texttt{idx\_fx}, \texttt{idx\_fy},
\texttt{idx\_fz}, \texttt{idx\_node} and \texttt{idx\_cell}, each taking $(i,j,k)$ and returning an offset into
the corresponding block. The six families share one flat index space, which keeps a field as a single array and
lets a source be described by a list of edge indices. The index helpers are also what the tests use to write a
stencil by hand and compare it with the assembled row, which is how the sign conventions of
\eqref{eq:identities} are pinned down.

The operators are assembled once, from coordinate triplets, and cached on the complex: $\Cmat$ of shape
$(n_{\mathrm{faces}},n_{\mathrm{edges}})$, $\Gmat$ of shape $(n_{\mathrm{edges}},n_{\mathrm{nodes}})$ and $\Smat$
of shape $(n_{\mathrm{cells}},n_{\mathrm{faces}})$. The two dual operators are not stored; they are the
transposes \texttt{curl\_dual()} $=\Cmat^{\!\top}$ and \texttt{divergence\_dual()} $=\Smat^{\!\top}$, exposed as
properties so that callers never index a transpose they have built themselves. This is the one place in the
package where it is easy to go wrong: \texttt{divergence()} builds one row per \emph{cell} with six face
contributions and outward signs, and therefore has shape $(n_{\mathrm{cells}},n_{\mathrm{faces}})$, while the
identities are often written with the divergence mapping faces to cells and the curl mapping edges to faces. Both
conventions are consistent, and a mismatch shows up immediately as the failure of \eqref{eq:identities} -- which
is exactly what the first tests assert, so the convention is checked on every run rather than remembered.

The stars are diagonal, so they are not stored as matrices at all. \texttt{material\_matrices} returns two
one-dimensional arrays, one entry per edge and one per face, and an application is an elementwise product. The
stepper caches their reciprocals once at construction, so the update \eqref{eq:leapfrog-h}--\eqref{eq:leapfrog-e}
costs two sparse products, two elementwise multiplications and one gather per step, with no allocation.
\texttt{MaterialMatrix} exists alongside the arrays for the eigensolver and for callers who prefer the linear
algebra to look like linear algebra; it holds the diagonal, its inverse and its solve, and supports
\texttt{@}. Material parameters may be given as a scalar, as a per-dof array, or as a vectorised callable of
$(x,y,z)$ sampled at edge or face centres -- the form \texttt{test\_energy} uses to build a two-layer medium
whose ratio must come out at exactly two.

The conductor constraint is a boolean mask. \texttt{pec\_edge\_mask} returns, for a choice of walls among the six
outer ones, a mask that is true exactly on the edges tangential to a conducting wall; the eigensolver uses it to
restrict the unknown space and the stepper uses it as an index, so that \eqref{eq:leapfrog-pec} is the single
assignment \texttt{e[mask] = 0.0} per step. On the demonstration grid it selects $620$ of the $1382$ edges, which
is the count the eigenvalue argument of Section~\ref{sec:nullspace} predicts without any geometric reasoning.

\subsection{A worked example}
\label{sec:api}

The following is a complete run, and is exactly what the package offers at its lowest level: no configuration
object, no solver class to inherit from, six lines from grid to field.

\begin{lstlisting}[style=python]
from emfit.grid import Grid
from emfit.operators import CellComplex
from emfit.materials import material_matrices
from emfit.evolve import FITStepper, RickerSource, cfl_dt

grid = Grid.uniform(40, 6, 6, lx=0.8, ly=0.12, lz=0.12)          # h = 0.02 m everywhere
cc   = CellComplex(grid)                                         # G, C, S and the PEC mask
m_eps, m_mu = material_matrices(grid, cc, eps_r=1.0)             # Hodge stars, one per dof
step = FITStepper(cc, m_eps, m_mu, dt=cfl_dt(grid), pec=cc.pec_edge_mask())
src  = RickerSource(cc.n_edges, idx=[cc.idx_ex(10, 3, 3)], fmax=1.5e9)

for _ in range(400):
    step.step(src(step.time))                                    # soft current on one edge
print(step.n, step.time, step.energy())                          # 400, 1.53e-8 s, 3.19e-9 J
\end{lstlisting}

The grid is the validation grid of Section~\ref{sec:energy} at $h=0.02$ m: $5404$ edges, $4836$ faces,
$\Delta t=3.81\times10^{-11}$ s from the bound \eqref{eq:cfl}, and $400$ steps in a few tens of milliseconds of wall-clock
time on a single core. Nothing is hidden in the six lines. The grid is built first because it carries the metric,
the complex from it because it carries the topology, the stars from both because they need both, and the stepper
is handed the reciprocal stars because that is all it uses. \texttt{step.step} takes the source values as an
argument rather than storing them, which makes a time-dependent excitation a closure and keeps the driver loop in
the caller's hands; \texttt{step.run(nsteps, source=..., probe=..., callback=...)} wraps the same loop for the
common case, and is what the demonstration uses.

Two consequences of the layout are visible from outside. Because a field is one array, a diagnostic can be a
function of that array alone, and \texttt{energy\_components()} returns the four terms of \eqref{eq:invariant}
separately rather than a single number -- which is how the wrong-sign and missing-term variants of
Table~\ref{tab:candidates} were measured. And because the constraint is a mask rather than a boundary condition
inside the solver, the absorbing condition of Section~\ref{sec:res-abc} was an addition to the same array and not a
change to the data layout: it updates edges in place and introduces no new unknowns. The layer of
Section~\ref{sec:res-pml} is as cheap, and for the same reason. It never enters the material tensor: it is a matched
lossy medium applied as one multiplication per field update on the arrays the stepper already touches.

\subsection{Testing}
\label{sec:testing}

The suite is seventy-one \texttt{pytest} tests in ten modules and takes $10$--$20$~s to run, of which the
cavity module accounts for roughly half. That imbalance is deliberate, and is worth stating plainly: the algebra
and the conservation laws are cheap to check, and a dense eigenvalue solve on a $14\times10\times7$ grid is not.

\begin{table}[H]
  \centering
  \caption{The ten test modules, what each one is allowed to claim, and how many tests it holds.}
  \label{tab:tests}
  \begin{tabular}{@{}l p{8.8cm} r@{}}
    \toprule
    Module & Claim under test & Tests \\
    \midrule
    \texttt{test\_operators} & $\Cmat\Gmat=0$, $\Smat\Cmat=0$, shapes, transposes, signs, metric independence & 15 \\
    \texttt{test\_energy}    & star values and units, exact conservation, CFL sharpness, wave speed & 6 \\
    \texttt{test\_abc}       & Mur absorption against the PEC control, $K\to0$ oracle, neighbour map & 5 \\
    \texttt{test\_pml}       & $\sigma_{\max}$ closed form, profile shape, multipliers, layer fit, below Mur & 7 \\
    \texttt{test\_periodic}  & seam identified with index zero, wrapped dual metric, dispersion & 9 \\
    \texttt{test\_cavity}    & spectrum against analytic values, null-space count, order of convergence & 9 \\
    \texttt{test\_scatter}   & staircase area, zero-radius control, forward-scattered energy & 3 \\
    \texttt{test\_nearfield} & DFT accumulator, Kirchhoff amplitude against a point source, RCS against Mie & 5 \\
    \texttt{test\_mie}       & truncation, losslessness, optical theorem, pinned pattern, large-$ka$ limit & 5 \\
    \texttt{test\_webapp}    & routing, form rendering, validation error, a real study and its figure & 7 \\
    \midrule
    \textbf{total} & & 71 \\
    \bottomrule
  \end{tabular}
\end{table}

The algebra module is the cheapest and the most useful. It checks that the two compositions vanish, that the
convenience views agree with the transposes, that every entry of the assembled operators is in
$\{-1,0,+1\}$, that a hand-written gradient stencil matches the assembled row, and that a jittered non-uniform
grid of the same size produces the same integer arrays -- the metric-independence claim of
Section~\ref{sec:modules}. These tests run in about a second and would catch any change to the numbering, which
is the kind of change that otherwise leaves the physics plausible and the answers wrong.

The time-domain module checks four things against closed forms rather than against stored output. The star
entries are compared with $A_{\tilde e}/L_e$ and $A_f/L_{\tilde f}$ computed by hand, and with a material ratio
that must come out at exactly two. The invariant \eqref{eq:invariant} is checked on random fields over $2000$
steps, with a tolerance of $10^{-6}$ against a measured drift of $1.3\times10^{-15}$, so the test asserts a
structural property and not a tuned number. The stability bound is checked on both sides: a step at $0.98$ of
\eqref{eq:cfl} stays bounded, and a step above the bound grows without limit. And the wave speed is measured
against the dispersion relation \eqref{eq:phase} on a grid ten cells per nominal wavelength, which is the
measurement reported in Section~\ref{sec:results}.

The cavity module is the expensive one and is the reason the suite is not instantaneous. A $14\times10\times7$
grid is large enough that a dense symmetric solve takes a few seconds -- about five in this container, and
plausibly more on a slower machine -- and small enough that its answer can be
compared with a fine-grid reference computed from the same code. What it establishes is the order of convergence,
the magnitude of the error on the two grids used in the report, the count of rejected static modes, and the
equality of the degenerate pairs -- all numbers quoted in Section~\ref{sec:eigen}. It also holds the one test
that asserts a limit: above the dense threshold the eigensolver must raise rather than silently attempt an
infeasible dense solve.

There are no golden files anywhere in the suite. Every expected value is either an integer count, a closed form,
or a number produced in the same run by an independent path, so a failing test names a broken law rather than a
changed output. The price of that discipline is the $28$--$42$~s of dense eigensolve; the benefit is that the
tolerances in the suite can be read as statements about the method rather than about the machine it last ran on.
